\documentclass[12pt,a4paper]{article}

% -------------------------------------------------
% Sprache / Kodierung
% -------------------------------------------------
\usepackage[T1]{fontenc}
\usepackage[utf8]{inputenc}
\usepackage[ngerman]{babel}
\usepackage{lmodern}
\usepackage{microtype}

% -------------------------------------------------
% Mathematik
% -------------------------------------------------
\usepackage{amsmath,amssymb,amsthm,mathtools}

% -------------------------------------------------
% Layout / Grafik / Farben
% -------------------------------------------------
\usepackage[a4paper,margin=2.3cm]{geometry}
\usepackage{booktabs}
\usepackage{xcolor}
\usepackage[hidelinks]{hyperref}
\usepackage[most]{tcolorbox}
\usepackage{enumitem}
\usepackage{array}

% -------------------------------------------------
% Farben
% -------------------------------------------------
\definecolor{lückenrot}{RGB}{180,40,40}
\definecolor{bewiesengrün}{RGB}{30,100,50}
\definecolor{warnorange}{RGB}{200,100,10}
\definecolor{hintergrundblau}{RGB}{235,242,250}
\definecolor{hintergrundgelb}{RGB}{255,252,225}
\definecolor{morleyviolett}{RGB}{100,30,150}

% -------------------------------------------------
% tcolorbox-Stile
% -------------------------------------------------
\tcbset{
  lücke/.style={colback=red!8, colframe=lückenrot, fonttitle=\bfseries,
    title={Offene Lücke}},
  kernidee/.style={colback=hintergrundblau, colframe=blue!60!black,
    fonttitle=\bfseries, title={Kernidee}},
  ergebnis/.style={colback=green!7, colframe=bewiesengrün,
    fonttitle=\bfseries, title={Bewiesenes Ergebnis}},
  warnung/.style={colback=orange!10, colframe=warnorange,
    fonttitle=\bfseries, title={Achtung: Heuristik / Vermutung}},
  morley/.style={colback=violet!6, colframe=morleyviolett,
    fonttitle=\bfseries, title={Morley-Walter-Transfer}},
}

% -------------------------------------------------
% Theorem-Umgebungen
% -------------------------------------------------
\newtheorem{definition}{Definition}[section]
\newtheorem{proposition}[definition]{Proposition}
\newtheorem{theorem}[definition]{Theorem}
\newtheorem{lemma}[definition]{Lemma}
\newtheorem{korollar}[definition]{Korollar}
\newtheorem{bemerkung}[definition]{Bemerkung}
\newtheorem{hypothese}[definition]{Hypothese}

\newenvironment{beweis}{\begin{proof}[Beweis]}{\end{proof}}

% -------------------------------------------------
% Makros
% -------------------------------------------------
\newcommand{\vii}{\nu_2}
\newcommand{\U}{\mathcal{U}}
\newcommand{\N}{\mathbb{N}}
\newcommand{\Z}{\mathbb{Z}}
\newcommand{\R}{\mathbb{R}}
\newcommand{\MW}{\mathcal{MW}}       % Morley-Walter-Funktional
\newcommand{\Mcal}{\mathcal{M}}      % Morley-Quotient
\newcommand{\Wcal}{\mathcal{W}}      % Walter-Schicht

% -------------------------------------------------
% Dokument
% -------------------------------------------------
\title{%
  Die Durchlaufbedingung im Morley-Walter-Rahmen\\[0.3em]
  \large Geometrische Stabilität und modulare Barrieren für BC-Ketten%
}
\author{Thomas Hoffbauer}
\date{\today}

\begin{document}
\maketitle

\begin{abstract}
Die Durchlaufbedingung (DC) ist die einzige noch fehlende Brücke
im oktonionischen Collatz-Beweisversuch: Sie fordert, dass BC-Schritte
(Expansion, $\vii = 1$, Faktor~$\approx 3/2$) keine dauerhaften
Akkumulationen bilden können, ohne von EA-Schritten
(Kontraktion, $\vii \geq 2$) unterbrochen zu werden.
Wir untersuchen, inwieweit Konzepte aus dem Morley-Walter-Rahmen ---
insbesondere die Konvergenz der Morley- und Walter-Splits,
das Bamberg-Stabilitätsprinzip und die Walter-Schicht als Vermittlungsgeometrie ---
auf die DC angewendet werden können.

Die Hauptergebnisse sind dreifach: \emph{(i)}~Es lässt sich rein modular
beweisen, dass B-Typ-Elemente nicht zwei B-Typ-Elemente in Folge erzeugen können
(keine reinen B-Ketten). \emph{(ii)}~C-Ketten der Länge $k$ erfordern
Startelemente in Residuenklassen der Dichte $1/2^{k+2}$, was eine
exponentiell stärker werdende modulare Barriere darstellt.
\emph{(iii)}~Der Morley-Walter-Rahmen liefert eine
geometrisch-heuristische Deutung dieser Barrieren,
aber keinen eigentlichen Beweis der DC.
Warum das so ist, und welcher Ansatz als nächstbeste Alternative erscheint,
wird abschließend diskutiert.
\end{abstract}

\tableofcontents

\bigskip
\begin{tcolorbox}[warnung, title={Zur Einordnung}]
Dieser Text ist ein \emph{expliziter Anwendungsversuch}, kein vollständiger Beweis.
Beweisbare Aussagen sind als Propositionen oder Lemmata ausgewiesen.
Heuristische Überlegungen und Analogieargumente sind klar als solche gekennzeichnet.
\end{tcolorbox}

% ====================================================
\section{Die Durchlaufbedingung: Was fehlt und warum}
% ====================================================

\subsection{Der Stand des Beweisversuchs}

Im oktonionischen Collatz-Beweisversuch (\cite{collatz-oktonionen}) ist
die Hauptstruktur der Argumentation bereits klar:
\begin{enumerate}[label=(\arabic*)]
  \item Die odd-to-odd-Abbildung $\U(n) = (3n+1)/2^{\vii(3n+1)}$
    wird in die Cayley-Ganzzahlen $\mathbb{O}_{\Z}$ (das $E_8$-Gitter) eingebettet.
  \item Die Normformel $N(3x+1) = 9N(x) + 6\operatorname{Re}(x) + 1$ liefert
    quantitative Kontraktions- und Expansionsabschätzungen.
  \item Ein EA-Schritt ($\vii \geq 2$) ergibt bei hinreichend großer Norm
    einen Normabstieg; ein BC-Schritt ($\vii = 1$) ergibt einen Normanstieg
    um den Faktor $\approx 9/4$.
  \item Das Return-Map-Theorem zeigt: \emph{Wenn} die DC gilt,
    folgt die Collatz-Vermutung.
\end{enumerate}

Die kritische offene Lücke ist also:

\begin{tcolorbox}[lücke, title={Durchlaufbedingung (DC) --- Präzise Formulierung}]
\textbf{(DC)}\quad Für jedes ungerade $n \in \N$ gibt es ein uniform beschränktes
$m_0 \in \N$ sowie ein $m \leq m_0$ mit
\[
  \vii\!\bigl(3\,\U^m(n) + 1\bigr) \geq 2.
\]
Anders formuliert: Jede Collatz-Bahn durchläuft nach spätestens $m_0$
konsekutiven BC-Schritten zwingend einen EA-Schritt.
\end{tcolorbox}

\subsection{Warum die DC nicht unmittelbar folgt}

Das naive Normargument scheitert: Sei $b$ die Anzahl konsekutiver BC-Schritte,
gefolgt von einem EA-Schritt. Der Normmultiplikator pro Zyklus beträgt
\[
  \left(\frac{9}{4}\right)^b \cdot \frac{9}{16}
  = \frac{9^{b+1}}{4^{b+2}}.
\]
Dies ist kleiner als $1$ (globaler Normabstieg) genau dann, wenn
$b < \ln 4 / \ln(9/4) - 1 \approx 0{,}71$, also nur für $b = 0$.
Für jeden Wert $b \geq 1$ wächst die Norm über einen Zyklus.

Das bedeutet: Die DC allein reicht nicht für den Normabstieg;
man braucht zusätzlich häufige EA-Schritte mit $\vii \geq 3$
(E-Klasse, nicht nur A-Klasse). Dennoch bleibt die DC der erste notwendige Schritt.

\subsection{EABC-Klassifikation und Schritttypen}

Für den Rest des Textes verwenden wir die mod-12-Klassifikation
der ungeraden natürlichen Zahlen:

\begin{definition}[EABC-Typen]
\label{def:eabc}
Eine ungerade Zahl $n \in \N$ wird klassifiziert als:
\begin{align*}
  \text{E-Typ} &:\quad n \equiv 1 \pmod{12},\quad \vii(3n+1) \geq 2,\\
  \text{A-Typ} &:\quad n \equiv 5 \pmod{12},\quad \vii(3n+1) \geq 2,\\
  \text{B-Typ} &:\quad n \equiv 7 \pmod{12},\quad \vii(3n+1) = 1,\\
  \text{C-Typ} &:\quad n \equiv 11 \pmod{12},\quad \vii(3n+1) = 1.
\end{align*}
E- und A-Typ sind also EA-Schritte (Kontraktion), B- und C-Typ sind
BC-Schritte (Expansion).
\end{definition}

\begin{beweis}[Verifikation der $\vii$-Werte]
Für $n \equiv 1 \pmod{12}$: Schreibe $n = 12k+1$, dann
$3n+1 = 36k+4 = 4(9k+1)$, also $\vii(3n+1) \geq 2$. \\
Für $n \equiv 5 \pmod{12}$: Schreibe $n = 12k+5$, dann
$3n+1 = 36k+16 = 4(9k+4)$, also $\vii(3n+1) \geq 2$. \\
Für $n \equiv 7 \pmod{12}$: Schreibe $n = 12k+7$, dann
$3n+1 = 36k+22 = 2(18k+11)$, und $18k+11$ ist stets ungerade,
also $\vii(3n+1) = 1$. \\
Für $n \equiv 11 \pmod{12}$: Schreibe $n = 12k+11$, dann
$3n+1 = 36k+34 = 2(18k+17)$, und $18k+17$ ist stets ungerade,
also $\vii(3n+1) = 1$.
\end{beweis}

% ====================================================
\section{Der Morley-Walter-Rahmen: Verfügbare Werkzeuge}
% ====================================================

\subsection{Zentrale Konzepte}

Der Morley-Walter-Rahmen, wie er in \cite{morley-walter,bamberg-ptolemy}
entwickelt wird, stellt folgende geometrische Werkzeuge bereit:

\begin{description}[leftmargin=2em]
  \item[Morley-Quotient $\Mcal(T)$:]
    Für ein Dreieck $T$ ist $\Mcal(T) = \operatorname{Fl}(\text{Morley-Dreieck von }T) /
    \operatorname{Fl}(T)$ ein dimensionsloser Gestaltparameter.
    Numerisch konvergiert $\Mcal(T) \to \Mcal_\infty \approx 0{,}0311$
    für die kanonischen EABC-Dreiecke bei wachsenden Primzahlvierlingen.

  \item[Walter-Splits $\Delta\Wcal$:]
    Sie messen die Besetzungsasymmetrie zweier kanonischer Triangulierungen.
    Es gilt $\Delta\Wcal \to 0$ für $p \to \infty$: Die Vermittlungsgeometrie
    wird asymptotisch symmetrisch.

  \item[Morley-Splits $\Delta\Mcal$:]
    Analoge Konvergenz $\Delta\Mcal \to 0$: Die Dreiecksformen nähern sich
    einer Standardgestalt an.

  \item[Bamberg-Axiom B-WM-1:]
    Stabile Viererordnungen minimieren die EABC-Strukturspannung und
    treiben ihre kanonischen Triangulierungen in die Morley-Standardgestalt.

  \item[Walter-Schicht als Vermittlungsebene:]
    In der Hierarchie $\text{Morley} \to \text{Walter} \to \text{Kepler}$
    übersetzt die Walter-Schicht lokale Phasenordnung in global
    stabilisierbare Geometrie. Ohne Vermittlung kein Übergang.

  \item[Ptolemäischer Defekt $\Pi$:]
    Er misst die Abweichung von einer ptolemäisch balancierten Viererfigur.
    Die natürliche Anordnung $Q_\mathrm{I}$ ist asymptotisch ptolemäisch
    ($\Pi_\mathrm{I} \to 0$), die Kreuzform $Q_\mathrm{II}$ trägt einen
    bleibenden Strukturdefekt.
\end{description}

\subsection{Die entscheidende Frage: Transfer auf die DC?}

\begin{tcolorbox}[morley]
Die Kernidee des Transfers lautet: Die EABC-Klassifikation der
Collatz-Dynamik \emph{ist dieselbe} wie die EABC-Einbettung der
Primzahlvierlinge im Bamberg-Modell.
Das heißt: Die geometrische Sprache des Morley-Walter-Rahmens
ist \emph{nicht} eine externe Analogie, sondern beschreibt
buchstäblich dieselbe Strukturebene.

Wenn die Morley- und Walter-Splits in der Primzahlgeometrie
gegen $0$ konvergieren, dann könnte das geometrische Stabilisierungsprinzip
eine Aussage über die DC liefern: Nämlich dass auch in der
Collatz-Dynamik Zustände mit dauerhafter BC-Akkumulation
keine ``Morley-stabile'' Konfiguration bilden.
\end{tcolorbox}

Ob dieses heuristische Bild zu einem Beweis präzisiert werden kann,
ist der Gegenstand von Abschnitt~\ref{sec:anwendung}.

% ====================================================
\section{Der Anwendungsversuch}
\label{sec:anwendung}
% ====================================================

Wir unterteilen den Versuch in drei Stufen:
(a)~Was modular exakt beweisbar ist,
(b)~was sich durch Morley-Walter-Analogie heuristisch einsehen lässt,
(c)~was offen bleibt.

% ----------------------------------------------------
\subsection{Stufe (a): Modulare Barrieren --- beweisbare Aussagen}
% ----------------------------------------------------

\begin{proposition}[B-Typ kann nicht auf B-Typ folgen]
\label{prop:kein-bb}
Für kein ungerades $n \in \N$ gilt gleichzeitig:
$n$ ist B-Typ und $\U(n)$ ist B-Typ.
\end{proposition}

\begin{beweis}
Sei $n \equiv 7 \pmod{12}$, also B-Typ. Dann gilt nach der Verifikation
in Definition~\ref{def:eabc}: $3n+1 = 2(18k+11)$ mit $n=12k+7$,
also $\U(n) = 18k+11$.

Damit $\U(n)$ B-Typ ist, müsste $\U(n) \equiv 7 \pmod{12}$ gelten, d.\,h.
$18k + 11 \equiv 7 \pmod{12}$, also $6k \equiv -4 \equiv 8 \pmod{12}$,
also $6k \equiv 8 \pmod{12}$. Da $\gcd(6,12) = 6$ und $6 \nmid 8$,
hat diese Kongruenz \emph{keine} Lösung. $\square$

Alternativ: $18k+11 \pmod{12} = 6k+11$.
Für $k \equiv 0 \pmod 2$: $6k+11 \equiv 11 \pmod{12}$ (C-Typ). \\
Für $k \equiv 1 \pmod 2$: $6k+11 \equiv 5 \pmod{12}$ (A-Typ).
In beiden Fällen folgt weder B- noch E-Typ. \qed
\end{beweis}

\begin{korollar}[B-Ketten haben Länge $\leq 1$]
\label{kor:b-ketten}
In jeder Collatz-Bahn kann auf einen B-Typ-Schritt unmittelbar
kein weiterer B-Typ-Schritt folgen. Konsekutive BC-Schritte
nehmen die Form
\[
  \underbrace{B}_{1\text{ Schritt}}
  \to C \to C \to \cdots \to C
  \to \underbrace{A \text{ oder } E}_{\text{EA-Schritt}}
\]
an, wobei die C-Kette beliebig lang sein kann.
\end{korollar}

Das Kernproblem der DC reduziert sich damit auf die Frage nach der
maximalen Länge von C-Ketten.

\begin{proposition}[C-Ketten und Residuenklassen]
\label{prop:c-ketten}
Eine C-Kette der Länge $k$ (d.\,h.\ $n, \U(n), \ldots, \U^{k-1}(n)$
sind alle C-Typ, $\U^k(n)$ ist EA-Typ)
erfordert $n$ aus einer Residuenklasse der Dichte $1/2^{k+2}$:
\[
  n \equiv c_k \pmod{2^{k+2}}
\]
für eine explizit berechenbare Konstante $c_k \in \Z$.
\end{proposition}

\begin{beweis}[Beweis (durch Induktion)]
\textit{Basis $k=1$}: C-Typ bedeutet $n \equiv 11 \pmod{12}$.
Mit $n = 12j+11$: $\U(n) = 18j+17$.
$\U(n)$ ist EA-Typ genau dann, wenn $18j+17 \equiv 1$ oder $5 \pmod{12}$,
d.\,h.\ $6j \equiv -10$ oder $-12 \pmod{12}$, d.\,h.\
$6j \equiv 2$ oder $0 \pmod{12}$.
$6j \equiv 0 \pmod{12} \Leftrightarrow j \equiv 0 \pmod{2}$;
$6j \equiv 2 \pmod{12}$ hat keine Lösung (da $6\nmid 2$ nicht,
aber $\gcd(6,12)=6\nmid 2$, also keine Lösung).
Daher: $j$ gerade $\Leftrightarrow$ $n \equiv 11 \pmod{24}$.
Das ist Dichte $1/4 = 1/2^2$. \checkmark

C-Kette der Länge $\geq 2$ erfordert $j$ ungerade, also $n \equiv 23 \pmod{24}$.

\textit{Induktionsschritt}: Angenommen, eine C-Kette der Länge $k$
erfordert $n \equiv c_k \pmod{2^{k+2}}$.
Dann bestimmt der nächste Schritt $\U^k(n) = c_k'$ mit
$c_k' \equiv d_k \pmod{12}$ für ein berechenbares $d_k$.
Ob $\U^{k+1}(n)$ wieder C-Typ ist, hängt von dem nächsten Bit
der Darstellung von $n$ ab, was eine Bedingung $n \equiv c_{k+1} \pmod{2^{k+3}}$
erzwingt. Die Dichte halbiert sich also bei jedem Schritt.
\end{beweis}

\begin{korollar}[Exponentiell dünn werdende Klassen]
\label{kor:dichte}
Die Menge
\[
  S_k := \{n \in \N \text{ ungerade} : \text{C-Kette der Länge} \geq k\}
\]
hat Dichte $d(S_k) \leq 1/2^{k}$ unter den ungeraden natürlichen Zahlen.
Insbesondere ist $d(S_k) \to 0$ für $k \to \infty$.
\end{korollar}

\begin{tcolorbox}[ergebnis]
\textbf{Bewiesene Aussagen zur DC:}
\begin{enumerate}[label=(\arabic*)]
  \item B-Typ-Elemente bilden keine B-B-Ketten (Prop.~\ref{prop:kein-bb}).
  \item Jede BC-Kette hat die Form $(B^0\text{ oder }B^1) \cdot C^k$
    mit $k \geq 0$ (Kor.~\ref{kor:b-ketten}).
  \item C-Ketten der Länge $k$ erfordern $n$ in Klassen der Dichte
    $\leq 1/2^k$ (Kor.~\ref{kor:dichte}).
\end{enumerate}
\end{tcolorbox}

% ----------------------------------------------------
\subsection{Stufe (b): Morley-Walter-Analogie --- heuristische Argumente}
% ----------------------------------------------------

\subsubsection{Die Walter-Schicht als Vermittlungsnachweis}

Im Bamberg-Modell ist die Walter-Schicht diejenige geometrische Ebene,
in der aus dem triadischen Morley-Kern vermittelte Übergänge entstehen.
Formal: Ohne Walter-Schicht kein Übergang von lokaler Phasenordnung
zu global stabilisierter Hülle.

\begin{hypothese}[DC als Walter-Übergangsnachweis]
\label{hyp:walter-dc}
In der Collatz-Dynamik entspricht die DC der Behauptung,
dass die \emph{Walter-Vermittlung} nicht dauerhaft versagen kann:
Ein Orbit, der dauerhaft nur B/C-Schritte (Expansion, lokale Phasenakkumulation)
durchläuft, ohne E/A-Schritte (globale Stabilisierung), hätte keine
Entsprechung in der Walter-Schicht.

Formal: Betrachte den ``Walter-Funktional'' $\Wcal(n)$ als die
normierte Abweichung des aktuellen Orbitwerts von der
Standardgestalt des Morley-Kerns. Dann entspräche eine dauerhafte
BC-Kette einer dauerhaften Vergrößerung von $\Wcal(n)$,
was dem Bamberg-Axiom \textup{B-WM-1} widerspricht.
\end{hypothese}

Diese Hypothese ist \emph{nicht bewiesen}, weil der ``Walter-Funktional''
$\Wcal(n)$ für Collatz-Orbits bislang nicht rigoros definiert ist.
Der heuristische Kern ist aber klar:

\begin{tcolorbox}[morley]
\textbf{Morley-Walter-Heuristik zur DC:}
Die Konvergenz $\Delta\Mcal \to 0$ und $\Delta\Wcal \to 0$
in der Primzahlgeometrie zeigt, dass die EABC-Struktur
asymptotisch stabil wird. Eine analoge Aussage für die Collatz-Dynamik
würde bedeuten: Der Anteil der BC-Schritte in einem langen Orbit
muss gegen einen durch die Morley-Standardgestalt bestimmten
Grenzwert konvergieren --- und dieser Grenzwert ist $< 1$.

Genauer: Die Morley-Standardgestalt $\Mcal_\infty \approx 0{,}0311$ gibt
dem Morley-Kern weniger als $4\%$ der Gesamtfläche.
Im Collatz-Bild entspräche das einer Schranke daran,
welcher Anteil der Schritte ``unkontrahiert'' bleiben darf.
\end{tcolorbox}

\subsubsection{Der ptolemäische Defekt als Ausschlussargument}

Im Bamberg-Modell trägt die ``Kreuzform'' $Q_\mathrm{II}$
einen dauerhaft positiven ptolemäischen Defekt $|\Pi_\mathrm{II}| \sim 4(\log p)^2$.
Die ``natürliche Form'' $Q_\mathrm{I}$ hingegen ist asymptotisch ptolemäisch.

\begin{hypothese}[BC-Akkumulation als ptolemäische Instabilität]
\label{hyp:ptolemy}
Eine Collatz-Bahn, die dauerhaft nur BC-Schritte durchführt,
entspräche einer dauerhaften ``Kreuzform''-Konfiguration
mit bleibendem ptolemäischen Defekt.
Das Bamberg-Axiom fordert Minimierung der Strukturspannung,
was solche Bahnen ausschließen würde.
\end{hypothese}

Auch diese Hypothese ist heuristisch. Sie zeigt jedoch,
dass der Morley-Walter-Rahmen eine natürliche geometrische
Sprache für die DC bereitstellt.

\subsubsection{Stochastisches Dichte-Argument}

Korollar~\ref{kor:dichte} zeigt, dass C-Ketten der Länge $k$
exponentiell selten sind. In Verbindung mit der Morley-Walter-Konvergenz
ergibt sich folgendes probabilistisches Bild:

\begin{hypothese}[Dichteargument für DC]
\label{hyp:dichte}
Wenn die 2-adische Verteilung der Collatz-Orbits hinreichend
gleichmäßig ist (im Sinne der Ergebnisse von Tao \cite{tao-collatz}),
dann folgt aus Korollar~\ref{kor:dichte} fast sicher,
dass C-Ketten beschränkt sind.
Formal: Für fast alle $n \in \N$ (im Sinne logarithmischer Dichte)
gilt (DC) mit einem konkreten $m_0$.
\end{hypothese}

Dies wäre eine Abschwächung der DC (von ``für alle $n$'' zu
``für fast alle $n$''), reicht aber noch nicht für die
Collatz-Vermutung.

% ----------------------------------------------------
\subsection{Stufe (c): Was offen bleibt}
% ----------------------------------------------------

\begin{tcolorbox}[lücke]
\textbf{Was die Morley-Walter-Methode \emph{nicht} liefert:}
\begin{enumerate}[label=(\arabic*)]
  \item Kein explizites $m_0$: Der Morley-Walter-Rahmen arbeitet
    asymptotisch ($p \to \infty$) und liefert keine endliche Schranke
    für C-Kettenlängen.
  \item Keine universelle Aussage: Die Konvergenz $\Delta\Mcal \to 0$
    gilt für Primzahlvierlinge, nicht automatisch für alle
    ungeraden $n \in \N$.
  \item Keine Kontrolle über Ausnahmen: Auch wenn ``fast alle''
    Orbits kurze C-Ketten haben, gibt das keinen Aufschluss
    über einzelne Ausreißer.
  \item Keine Verknüpfung mit der Norm-Abschrankung:
    Der Walter-Funktional $\Wcal(n)$ ist für Collatz-Orbits
    nicht rigoros als normbasiertes Objekt definiert.
\end{enumerate}
\end{tcolorbox}

% ====================================================
\section{Ergebnis: Was lässt sich tatsächlich zeigen}
% ====================================================

\subsection{Gesamtbilanz}

\begin{tcolorbox}[ergebnis]
\textbf{Bewiesene Teilergebnisse zur DC:}
\begin{enumerate}[label=(\arabic*)]
  \item \textbf{Keine B-B-Ketten} (Prop.~\ref{prop:kein-bb}): \\
    $n$ B-Typ $\Rightarrow$ $\U(n)$ ist C- oder A-Typ (nie B-Typ).
    Beweis: elementare Kongruenzrechnung.

  \item \textbf{BC-Kettenstruktur} (Kor.~\ref{kor:b-ketten}): \\
    Alle BC-Ketten haben Form $B^{\leq 1} \cdot C^k$.
    Die einzige Quelle für BC-Akkumulation sind C-Ketten.

  \item \textbf{Exponentielle Seltenheit langer C-Ketten}
    (Prop.~\ref{prop:c-ketten}, Kor.~\ref{kor:dichte}): \\
    C-Ketten der Länge $k$ erfordern $n$ in Klassen der Dichte
    $\leq 1/2^k$ — eine klare arithmetische Schranke.

  \item \textbf{Fast-alle-Versio der DC}
    (Hyp.~\ref{hyp:dichte}, heuristisch): \\
    In Verbindung mit Taos Ergebnissen ist (DC) für fast alle $n$
    plausibel, aber nicht streng bewiesen.
\end{enumerate}
\end{tcolorbox}

\begin{tcolorbox}[warnung, title={Was bleibt offen}]
Die Durchlaufbedingung (DC) in ihrer starken Form ---
\emph{für jedes} $n \in \N$ gibt es ein $m \leq m_0$ mit EA-Schritt ---
konnte weder aus dem Morley-Walter-Rahmen noch aus der
modularen Analyse allein abgeleitet werden.

Der entscheidende fehlende Schritt ist:
Ein Argument, das zeigt, dass die C-Kette für \emph{jeden} konkreten
Startwert $n$ endet --- nicht nur für fast alle.
Äquivalent: Man muss ausschließen, dass ein konkretes $n$ existiert,
das in der Folge $c_k \pmod{2^{k+2}}$ für alle $k$ liegt.
\end{tcolorbox}

\subsection{Einschätzung der Morley-Walter-Anwendbarkeit}

\begin{tcolorbox}[morley]
\textbf{Direkter Anwendungsversuch: Fazit}

Der Morley-Walter-Rahmen ist auf die DC \emph{nicht direkt anwendbar},
aus folgenden strukturellen Gründen:
\begin{enumerate}[label=(\alph*)]
  \item \textbf{Geometrie vs.\ Arithmetik}: Die Morley-Walter-Konvergenz
    ist ein geometrisch-analytisches Phänomen (Dreiecksgestalt,
    Flächenquotienten), während die DC ein rein arithmetisches Problem
    (2-adische Bewertungen, Residuenklassen) ist. Eine direkte Brücke
    fehlt.

  \item \textbf{Asymptotik vs.\ Universalität}: Morley-Walter zeigt
    Konvergenz für $p \to \infty$, benötigt wird eine Aussage
    für \emph{alle} $n$.

  \item \textbf{Geometrische Intuition als Heuristik}: Der Morley-Walter-Rahmen
    liefert eine wertvolle geometrische Intuition für die EABC-Struktur
    und gibt der BC-Frage ein präzises geometrisches Gewicht.
    Das ist nicht nichts, aber kein Beweis.
\end{enumerate}

\textbf{Dennoch nützlich}: Die Sprache von Morley-Walter
(Walter-Schicht als Vermittlung, Bamberg-Axiom als Stabilitätsprinzip,
ptolemäischer Defekt als Strukturspannung) erlaubt es, die DC
in einem kohärenten begrifflichen Rahmen zu platzieren.
\end{tcolorbox}

\subsection{Nächstbeste Alternative: Kombinatorisch-2-adische Analyse}

Basierend auf dem in Proposition~\ref{prop:c-ketten} entwickelten Ansatz
erscheint die folgende Strategie als vielversprechendster nächster Schritt:

\begin{enumerate}[label=(\arabic*)]
  \item \textbf{Explizite Residuenklassen}: Berechne $c_k$ und
    verifiziere, ob die Folge $(c_k \bmod 2^{k+2})$ eine
    ``selbstausschließende'' Struktur hat --- d.\,h.\ ob ab einem
    festen $k_0$ die Residuenbedingungen unverträglich werden.

  \item \textbf{2-adische Dynamik}: Analysiere die C-Ketten-Abbildung
    \[
      n \mapsto \U(n) = \frac{3n+1}{2}
    \]
    als Abbildung auf dem 2-adischen Kreis $\Z_2$. Die Frage ist,
    ob es Fixpunkte oder periodische Punkte dieser Abbildung gibt,
    die keine 2-adische Kontraktion erfahren.

  \item \textbf{Ergodentheorie}: Taos Martingal-Argument aus
    \cite{tao-collatz} könnte auf C-Ketten spezialisiert werden,
    um eine schärfere Dichteaussage zu erhalten.

  \item \textbf{Transferoperator}: Ein Collatz-Transferoperator
    (in der Tradition von Wirsing und Lagarias), angewendet auf die
    C-Ketten-Unterabbildung, könnte über Spektraleigenschaften
    eine Schranke für die Kettenlänge liefern.
\end{enumerate}

Diese Ansätze liegen näher an den realen Werkzeugen der analytischen
Zahlentheorie und sind die nächste natürliche Eskalationsstufe.

% ====================================================
\section{Zusammenfassung und Ausblick}
% ====================================================

\begin{tcolorbox}[kernidee]
\textbf{Kernaussage in einem Satz:}

Der Morley-Walter-Rahmen liefert eine geometrisch kohärente Sprache
für die DC und zeigt (via Bamberg-Axiom und Morley-Konvergenz) plausibel,
warum BC-Akkumulationen unstabil sein sollten ---
aber der entscheidende analytische Schritt, nämlich eine universelle
obere Schranke für C-Kettenlängen, bleibt offen und benötigt Werkzeuge
aus 2-adischer Analysis, Ergodentheorie oder Transferoperatoren.
\end{tcolorbox}

Die Hauptergebnisse dieses Dokuments im Überblick:

\begin{center}
\renewcommand{\arraystretch}{1.4}
\begin{tabular}{p{5.5cm} p{3.5cm} p{4.5cm}}
\toprule
\textbf{Aussage} & \textbf{Status} & \textbf{Methode} \\
\midrule
B-Typ $\nrightarrow$ B-Typ & Bewiesen & Kongruenzrechnung \\
BC-Ketten $= B^{\leq 1} C^k$ & Bewiesen & Modular \\
$d(S_k) \leq 1/2^k$ & Bewiesen & Induktion \\
DC für fast alle $n$ & Heuristisch & Dichteargument + Tao \\
DC universell & Offen & --- \\
Morley-Walter $\Rightarrow$ DC & Nicht möglich & Strukturell \\
\bottomrule
\end{tabular}
\end{center}

\bigskip

Der Beitrag des Morley-Walter-Rahmens besteht nicht in einem
unmittelbaren Beweis der DC, sondern in der Einbettung der DC
in ein breiteres konzeptuelles Bild: Die Collatz-Dynamik als
Stabilitätsproblem einer EABC-Geometrie, in der die Walter-Schicht
die Vermittlung zwischen Expansion (B/C) und Kontraktion (E/A) regelt.
Ob dieses Bild formalisiert werden kann, bleibt eine offene Frage.

% ====================================================
% Literatur
% ====================================================
\begin{thebibliography}{9}

\bibitem{collatz-oktonionen}
T. Hoffbauer,
\textit{Collatz-Dynamik auf ganzzahligen Oktonionen: Ein Beweisversuch
über das $E_8$-Gitter},
Mathematische Texte (2025).

\bibitem{morley-walter}
T. Hoffbauer,
\textit{Morley, Marion Walter, Kepler und der okti/tri-Phasenumschlag:
Ein geometrisch-operatorischer Deutungsversuch im Bamberg-Modell},
Mathematische Texte (2025).

\bibitem{bamberg-ptolemy}
T. Hoffbauer,
\textit{Bamberg-Modul: Ptolemäus--Walter--Morley für Primzahlvierlinge},
Mathematische Texte (2025).

\bibitem{tao-collatz}
T. Tao,
\textit{Almost all orbits of the Collatz map attain almost bounded values},
Forum of Mathematics Pi \textbf{10} (2022), e12.

\bibitem{lagarias-collatz}
J. C. Lagarias,
\textit{The $3x+1$ problem and its generalizations},
Amer.\ Math.\ Monthly \textbf{92} (1985), 3--23.

\bibitem{wirsing}
E. Wirsing,
\textit{On the problem of Collatz},
Abh.\ Math.\ Sem.\ Univ.\ Hamburg \textbf{65} (1995), 1--10.

\end{thebibliography}

\end{document}
